\documentclass[10pt, letterpaper]{article}

% Packages:
\usepackage[
    ignoreheadfoot, % set margins without considering header and footer
    top=2 cm, % seperation between body and page edge from the top
    bottom=2 cm, % seperation between body and page edge from the bottom
    left=2 cm, % seperation between body and page edge from the left
    right=2 cm, % seperation between body and page edge from the right
    footskip=1.0 cm, % seperation between body and footer
    % showframe % for debugging
]{geometry} % for adjusting page geometry
\usepackage{titlesec} % for customizing section titles
\usepackage{tabularx} % for making tables with fixed width columns
\usepackage{array} % tabularx requires this
\usepackage[dvipsnames]{xcolor} % for coloring text
\definecolor{primaryColor}{RGB}{0, 0, 0} % define primary color
\usepackage{enumitem} % for customizing lists
% \usepackage{fontawesome5} % for using icons -- removed: unused, and it crashes
% the local Tectonic build. Re-add on Overleaf if you ever want icons.
\usepackage{amsmath} % for math
\usepackage[
    pdftitle={Chirag Lakhani's CV},
    pdfauthor={Chirag Lakhani},
    pdfcreator={LaTeX with RenderCV},
    colorlinks=true,
    urlcolor=primaryColor
]{hyperref} % for links, metadata and bookmarks
\usepackage[pscoord]{eso-pic} % for floating text on the page
\usepackage{calc} % for calculating lengths
\usepackage{bookmark} % for bookmarks
\usepackage{lastpage} % for getting the total number of pages
\usepackage{changepage} % for one column entries (adjustwidth environment)
\usepackage{paracol} % for two and three column entries
\usepackage{ifthen} % for conditional statements
\usepackage{needspace} % for avoiding page brake right after the section title
\usepackage{iftex} % check if engine is pdflatex, xetex or luatex



% Ensure that generate pdf is machine readable/ATS parsable:
\ifPDFTeX
    \input{glyphtounicode}
    \pdfgentounicode=1
    \usepackage[T1]{fontenc}
    \usepackage[utf8]{inputenc}
    \usepackage{lmodern}
\fi

\usepackage{charter}


\usepackage[style=numeric,sorting=ydnt,maxbibnames=99,defernumbers=true]{biblatex}
\addbibresource{chirag_citations.bib} % your .bib file name





% Some settings:
\raggedright
\AtBeginEnvironment{adjustwidth}{\partopsep0pt} % remove space before adjustwidth environment
\pagestyle{empty} % no header or footer
\setcounter{secnumdepth}{0} % no section numbering
\setlength{\parindent}{0pt} % no indentation
\setlength{\topskip}{0pt} % no top skip
\setlength{\columnsep}{0.15cm} % set column seperation
\pagenumbering{gobble} % no page numbering

\titleformat{\section}{\needspace{4\baselineskip}\bfseries\large}{}{0pt}{}[\vspace{1pt}\titlerule]

\titlespacing{\section}{
    % left space:
    -1pt
}{
    % top space:
    0.3 cm
}{
    % bottom space:
    0.2 cm
} % section title spacing

% Subsection headings (used inside Presentations):
\titleformat{\subsection}{\needspace{3\baselineskip}\bfseries\normalsize}{}{0pt}{}
\titlespacing{\subsection}{-1pt}{0.25 cm}{0.15 cm}

\renewcommand\labelitemi{$\vcenter{\hbox{\small$\bullet$}}$} % custom bullet points
\newenvironment{highlights}{
    \begin{itemize}[
        topsep=0.10 cm,
        parsep=0.10 cm,
        partopsep=0pt,
        itemsep=0pt,
        leftmargin=0 cm + 10pt
    ]
}{
    \end{itemize}
} % new environment for highlights


\newenvironment{highlightsforbulletentries}{
    \begin{itemize}[
        topsep=0.10 cm,
        parsep=0.10 cm,
        partopsep=0pt,
        itemsep=0pt,
        leftmargin=10pt
    ]
}{
    \end{itemize}
} % new environment for highlights for bullet entries

\newenvironment{onecolentry}{
    \begin{adjustwidth}{
        0 cm + 0.00001 cm
    }{
        0 cm + 0.00001 cm
    }
}{
    \end{adjustwidth}
} % new environment for one column entries

\newenvironment{twocolentry}[2][]{
    \onecolentry
    \def\secondColumn{#2}
    \setcolumnwidth{\fill, 4.5 cm}
    \begin{paracol}{2}
}{
    \switchcolumn \raggedleft \secondColumn
    \end{paracol}
    \endonecolentry
} % new environment for two column entries

\newenvironment{threecolentry}[3][]{
    \onecolentry
    \def\thirdColumn{#3}
    \setcolumnwidth{, \fill, 4.5 cm}
    \begin{paracol}{3}
    {\raggedright #2} \switchcolumn
}{
    \switchcolumn \raggedleft \thirdColumn
    \end{paracol}
    \endonecolentry
} % new environment for three column entries

\newenvironment{header}{
    \setlength{\topsep}{0pt}\par\kern\topsep\centering\linespread{1.5}
}{
    \par\kern\topsep
} % new environment for the header

\newcommand{\placelastupdatedtext}{% \placetextbox{<horizontal pos>}{<vertical pos>}{<stuff>}
  \AddToShipoutPictureFG*{% Add <stuff> to current page foreground
    \put(
        \LenToUnit{\paperwidth-2 cm-0 cm+0.05cm},
        \LenToUnit{\paperheight-1.0 cm}
    ){\vtop{{\null}\makebox[0pt][c]{
        \small\color{gray}\textit{Last updated in September 2026}\hspace{\widthof{Last updated in September 2026}}
    }}}%
  }%
}%

% save the original href command in a new command:
\let\hrefWithoutArrow\href

% new command for external links:


% Name commands: bold in every bibliography entry. The First/Second superscript
% variants are retired (author-position notes aren't standard on a CV) but kept
% as aliases so older .bib entries still compile.
\newcommand{\myname}[1]{\textbf{#1}}
\newcommand{\mynameFirst}[1]{\textbf{#1}}
\newcommand{\mynameSecond}[1]{\textbf{#1}}


\begin{document}
    \newcommand{\AND}{\unskip
        \cleaders\copy\ANDbox\hskip\wd\ANDbox
        \ignorespaces
    }
    \newsavebox\ANDbox
    \sbox\ANDbox{$|$}

    \begin{header}
        \fontsize{25 pt}{25 pt}\selectfont Chirag Lakhani

        \vspace{5 pt}

        \normalsize
        \mbox{New York, NY}%
        \kern 5.0 pt%
        \AND%
        \kern 5.0 pt%
        \mbox{\hrefWithoutArrow{https://orcid.org/0000-0002-8431-3690}{orcid.org/0000-0002-8431-3690}}%
        \kern 5.0 pt%
        \AND%
        \kern 5.0 pt%
        \mbox{\hrefWithoutArrow{}{clakhani (at) nygenome (dot) org}}%
        \kern 5.0 pt%
        \AND%
        \kern 5.0 pt%
        \mbox{\hrefWithoutArrow{https://cmlakhan.github.io/}{cmlakhan.github.io}}%
        \kern 5.0 pt%
        \AND%
        \kern 5.0 pt%
        \mbox{\hrefWithoutArrow{https://www.linkedin.com/in/chiraglakhani/}{linkedin.com/in/chiraglakhani/}}%
        \kern 5.0 pt%
        \AND%
        \kern 5.0 pt%
        \mbox{\hrefWithoutArrow{https://github.com/cmlakhan}{github.com/cmlakhan}}%
    \end{header}

    \vspace{5 pt - 0.3 cm}




\section{Research Summary}

    \begin{onecolentry}
I develop statistical and deep learning methods for predicting the functional effects of regulatory genetic variants, and apply them to identify causal variants and uncover missing regulation in Alzheimer's disease. My work integrates genomic deep learning models with functional genomics data to understand cell-type-specific gene regulation in neurodegenerative disease.

\end{onecolentry}






\section{Positions Held}

        \begin{twocolentry}{
            01/2021 – Present
        }
            \textbf{Senior Research Scientist}, New York Genome Center
            \end{twocolentry}

        \vspace{0.10 cm}
        \begin{onecolentry}
            \begin{highlights}
                \item Develop deep learning methods for predicting regulatory variant effects and apply them to fine-mapping and heritability analysis in Alzheimer's disease, in collaboration with David Knowles (Columbia University) and Towfique Raj (Icahn School of Medicine at Mount Sinai).
                \item Lead author on two methods papers integrating genomic deep learning models with functional genomics: functionally informed fine-mapping of causal AD variants, and cell-type-resolved eQTL prediction (scEEMS) explaining AD heritability.
                \item Collaborate with Gao Wang (Columbia University) on molecular QTL prediction within the ADSP Functional Genomics Consortium (FunGen-AD); funded as a researcher on NIH/NIA R01~AG086467 (\textit{FunGen-xQTL}, PI: Gao Wang) and previously on NIH/NIA U01~AG068880 (\textit{Learning the Regulatory Code of Alzheimer's Disease Genomes}).
                \item Present regularly to the ADSP AI/ML Working Group and its External Advisory Board.
            \end{highlights}
        \end{onecolentry}

        \vspace{0.20 cm}


        \begin{twocolentry}{
            03/2015 – 12/2020
        }
            \textbf{Postdoctoral Research Fellow}, Harvard Medical School - Department of Biomedical Informatics
            \end{twocolentry}

        \vspace{0.10 cm}
        \begin{onecolentry}
            \begin{highlights}
               \item Advisor: Chirag J. Patel (Harvard Medical School)
                \item Repurposed a health insurance claims dataset of 45 million individuals to conduct the largest twin study in the United States (56,396 twin pairs), estimating genetic and environmental contributions to 560 phenotypes. Published in \textit{Nature Genetics}; press coverage in the Washington Post, CNN, and The Verge.
                \item Developed the Exposome Data Warehouse, a database linking publicly available environmental data---socioeconomic, pollution, and climate---to electronic health records.
                \item NIH T32 Training Fellow, 2018--2020 (co-mentored by James Meigs, Massachusetts General Hospital).
            \end{highlights}
        \end{onecolentry}

        \vspace{0.20 cm}


        \begin{twocolentry}{
            04/2014 – 03/2015
        }
            \textbf{Data Scientist}, HelloWallet, Inc
            \end{twocolentry}

        \vspace{0.10 cm}
        \begin{onecolentry}
            \begin{highlights}
                \item Designed and deployed a machine learning transaction classification system in a production consumer application; led migration of analytics infrastructure to a distributed computing environment.
            \end{highlights}
        \end{onecolentry}


        \vspace{0.20 cm}


        \begin{twocolentry}{
            08/2013 – 04/2014
        }
            \textbf{Data Scientist}, Zaloni, Inc
            \end{twocolentry}

        \vspace{0.10 cm}
        \begin{onecolentry}
            \begin{highlights}
                \item Built the company's data science practice, delivering large-scale machine learning solutions for enterprise clients. Developed and taught an Introduction to Data Science course at several Fortune 500 companies.
            \end{highlights}
        \end{onecolentry}

        \vspace{0.20 cm}



        \begin{twocolentry}{
            01/2012 – 08/2013
        }
            \textbf{Visiting Researcher (unaffiliated)}, SAMSI and Duke University
            \end{twocolentry}

        \vspace{0.10 cm}
        \begin{onecolentry}
            \begin{highlights}
                \item Participated in SAMSI workshops and Mauro Maggioni's research group on geometric multi-resolution analysis and scalable manifold learning.
            \end{highlights}
        \end{onecolentry}


        \vspace{0.20 cm}



        \begin{twocolentry}{
            08/2011 – 12/2011
        }
            \textbf{Postdoctoral Research Associate}, North Carolina State University
            \end{twocolentry}

        \vspace{0.10 cm}
        \begin{onecolentry}
            \begin{highlights}
                \item Advisor: Hamid Krim (North Carolina State University)
                \item Applied computational topology to data mining: persistent cohomology for nonlinear dimensionality reduction, and Morse theory for community detection in networks.
            \end{highlights}
        \end{onecolentry}


        \vspace{0.20 cm}








\section{Education}




        \begin{twocolentry}{
            12/2010
        }
            \textbf{North Carolina State University}, PhD in Mathematics\end{twocolentry}

        \vspace{0.10 cm}
        \begin{onecolentry}
            \begin{highlights}
                \item Thesis: The GIT Compactification of Quintic Threefolds
                \item Advisor: Amassa Fauntleroy
            \end{highlights}
        \end{onecolentry}


        \vspace{0.20 cm}




        \begin{twocolentry}{
            12/2003
        }
            \textbf{North Carolina State University}, B.S. in Mathematics \end{twocolentry}

        \vspace{0.10 cm}
        \begin{onecolentry}
            \begin{highlights}
                \item Minor: Physics
                \item Magna Cum Laude
            \end{highlights}
        \end{onecolentry}

        \vspace{0.20 cm}

        \begin{onecolentry}
            \textbf{Additional Training:} Introduction to Statistical Genetics (BST~227, Martin Aryee, Fall 2016), Molecular Biology for Epidemiologists (EPI~249, Immaculata De Vivo, Fall 2016), and Advanced Population and Medical Genetics (EPI~511, Alkes Price, Spring 2017), Harvard T.H. Chan School of Public Health; 22nd Summer Institute in Statistical Genetics, University of Washington, workshop with Peter Visscher (July 2017).
        \end{onecolentry}





\section{Publications}
     \nocite{*} % This tells biblatex to include ALL entries from your .bib file
\printbibliography[keyword=preprint,heading=subbibliography,title={Preprints Under Review}]
\printbibliography[keyword=published,heading=subbibliography,title={Peer-Reviewed Publications}]
\printbibliography[keyword=earlier,heading=subbibliography,title={Earlier Work (Doctoral Training)}]



\section{Presentations}

\subsection{Conference Presentations}

    \begin{samepage}
        \begin{twocolentry}{
            August 2026
        }
            \textbf{Machine Learning-Based Prediction of Cell-type Resolved Brain eQTLs Enhances Discovery of Variants Explaining Alzheimer's Disease Heritability} (Talk)
        \end{twocolentry}

        \vspace{0.10 cm}

        \begin{onecolentry}
            NIH Alzheimer's Disease Sequencing Project Annual Meeting, Washington, DC
        \end{onecolentry}
    \end{samepage}

    \vspace{0.20 cm}

    \begin{samepage}
        \begin{twocolentry}{
            May 2026
        }
            \textbf{Integration of Deep Learning Annotations with Functional Genomics Improves Identification of Causal Alzheimer's Disease Variants} (Poster)
        \end{twocolentry}

        \vspace{0.10 cm}

        \begin{onecolentry}
            90th Cold Spring Harbor Symposium on Quantitative Biology: AI in Biology, Cold Spring Harbor, NY
        \end{onecolentry}
    \end{samepage}

    \vspace{0.20 cm}

    \begin{samepage}
        \begin{twocolentry}{
            July 2025
        }
            \textbf{Integration of Deep Learning Annotations with Functional Genomics Improves Identification of Causal Alzheimer's Disease Variants} (Poster)
        \end{twocolentry}

        \vspace{0.10 cm}

        \begin{onecolentry}
            Gordon Research Conference on Alzheimer's Disease, Ventura, CA
        \end{onecolentry}
    \end{samepage}

    \vspace{0.20 cm}

    \begin{samepage}
        \begin{twocolentry}{
            May 2025
        }
            \textbf{MLxQTL: Machine Learning Based Prediction of Single Cell eQTLs in Neuronal Cell Types} (Talk)
        \end{twocolentry}

        \vspace{0.10 cm}

        \begin{onecolentry}
            NIH Alzheimer's Disease Sequencing Project xQTL Workshop, San Francisco, CA
        \end{onecolentry}
    \end{samepage}

    \vspace{0.20 cm}

    \begin{samepage}
        \begin{twocolentry}{
            May 2024
        }
            \textbf{Localization of Polygenic Signal in Alzheimer's Disease through the Integration of Cell-Type Functional Annotations and Deep Learning Models} (Poster)
        \end{twocolentry}

        \vspace{0.10 cm}

        \begin{onecolentry}
            Biology of Genomes, Cold Spring Harbor, NY
        \end{onecolentry}
    \end{samepage}

    \vspace{0.20 cm}

    \begin{samepage}
        \begin{twocolentry}{
            March 2024
        }
            \textbf{Deep Learning Approaches for Functional Genomic Variant Prediction using xQTL Data} (Talk)
        \end{twocolentry}

        \vspace{0.10 cm}

        \begin{onecolentry}
            NIH Alzheimer's Disease Sequencing Project xQTL Workshop, New York, NY
        \end{onecolentry}
    \end{samepage}

    \vspace{0.20 cm}

    \begin{samepage}
        \begin{twocolentry}{
            March 2024
        }
            \textbf{Localization of Polygenic Signal in Alzheimer's Disease through the Integration of Cell-Type Functional Annotations and Deep Learning Models} (Poster)
        \end{twocolentry}

        \vspace{0.10 cm}

        \begin{onecolentry}
            Broad Institute Variant-to-Function Meeting, Cambridge, MA
        \end{onecolentry}
    \end{samepage}

    \vspace{0.20 cm}

    \begin{samepage}
        \begin{twocolentry}{
            October 2023
        }
            \textbf{Localization of Polygenic Signal in Alzheimer's Disease through the Integration of Cell-Type Functional Annotations and Deep Learning Models} (Poster)
        \end{twocolentry}

        \vspace{0.10 cm}

        \begin{onecolentry}
            American Society of Human Genetics, Washington, DC
        \end{onecolentry}
    \end{samepage}

    \vspace{0.20 cm}

    \begin{samepage}
        \begin{twocolentry}{
            September 2019
        }
            \textbf{Repurposing Large Health Insurance Claims Data and Publicly Available Exposomic Data to Estimate Genetic and Environmental Contributions in 560 Phenotypes} (Talk)
        \end{twocolentry}

        \vspace{0.10 cm}

        \begin{onecolentry}
            Chan-Zuckerberg Initiative Workshop, CU Anschutz, Aurora, CO
        \end{onecolentry}
    \end{samepage}

    \vspace{0.20 cm}

    \begin{samepage}
        \begin{twocolentry}{
            March 2017
        }
            \textbf{Building a Search Engine to Find Environmental and Phenotypic Factors Associated with Health and Disease} (Talk)
        \end{twocolentry}

        \vspace{0.10 cm}

        \begin{onecolentry}
            NSF Big Data Spokes Meeting, Washington, DC
        \end{onecolentry}
    \end{samepage}

    \vspace{0.20 cm}

    \begin{samepage}
        \begin{twocolentry}{
            October 2015
        }
            \textbf{Systematic and Large-Scale Investigation of Twin and Sibling Concordance of 1723 Traits in a Nationally Representative Health Claims Cohort} (Platform Talk)
        \end{twocolentry}

        \vspace{0.10 cm}

        \begin{onecolentry}
            American Society of Human Genetics, Baltimore, MD
        \end{onecolentry}
    \end{samepage}

    \vspace{0.20 cm}

\subsection{Consortium Presentations}

    \begin{samepage}
        \begin{twocolentry}{
            October 2024
        }
            \textbf{Learning the Regulatory Code of Alzheimer's Disease Genomes} (Talk)
        \end{twocolentry}

        \vspace{0.10 cm}

        \begin{onecolentry}
            NIH Alzheimer's Disease Sequencing Project AI/ML External Advisory Board Meeting, Bethesda, MD
        \end{onecolentry}
    \end{samepage}

    \vspace{0.20 cm}

    \begin{samepage}
        \begin{twocolentry}{
            August 2024
        }
            \textbf{Learning the Regulatory Code of Alzheimer's Disease Genomes} (Talk)
        \end{twocolentry}

        \vspace{0.10 cm}

        \begin{onecolentry}
            NIH Alzheimer's Disease Sequencing Project AI/ML Working Group, Online
        \end{onecolentry}
    \end{samepage}

    \vspace{0.20 cm}

    \begin{samepage}
        \begin{twocolentry}{
            March 2023
        }
            \textbf{Deep Learning Analysis of Genetic Variants in PICALM/EED Locus} (Talk)
        \end{twocolentry}

        \vspace{0.10 cm}

        \begin{onecolentry}
            NIH Alzheimer's Disease Sequencing Project Functional Genomics Meeting, Online
        \end{onecolentry}
    \end{samepage}

    \vspace{0.20 cm}

\subsection{Earlier Talks}

    \begin{samepage}
        \begin{twocolentry}{
            February 2016
        }
            \textbf{Building an Exposome API}
        \end{twocolentry}

        \vspace{0.10 cm}

        \begin{onecolentry}
            Environmental Statistics Seminar, Harvard T.H. Chan School of Public Health, Boston, MA
        \end{onecolentry}
    \end{samepage}

    \vspace{0.20 cm}

    \begin{samepage}
        \begin{twocolentry}{
            November 2012
        }
            \textbf{Approximation of Kernel Matrices}
        \end{twocolentry}

        \vspace{0.10 cm}

        \begin{onecolentry}
            SAMSI Working Group, SAMSI, Research Triangle Park, NC
        \end{onecolentry}
    \end{samepage}

    \vspace{0.20 cm}

    \begin{samepage}
        \begin{twocolentry}{
            October 2011
        }
            \textbf{Tensor Decompositions in Signal Processing}
        \end{twocolentry}

        \vspace{0.10 cm}

        \begin{onecolentry}
            Tensor Seminar, North Carolina State University, Raleigh, NC
        \end{onecolentry}
    \end{samepage}





\section{Mentoring}

    \begin{onecolentry}
        \textit{Students co-supervised at the New York Genome Center.}
    \end{onecolentry}

    \vspace{0.10 cm}

\subsection{Graduate Students}

    \begin{onecolentry}
        \textbf{Giacomo Cavalca}
    \end{onecolentry}

    \vspace{0.10 cm}
    \begin{onecolentry}
        \begin{highlights}
            \item Machine learning-based prediction of cell-type-resolved brain eQTLs; co-author on the scEEMS preprint.
        \end{highlights}
    \end{onecolentry}

    \vspace{0.20 cm}

\subsection{Undergraduate Students}

    \begin{onecolentry}
        \textbf{Rohan Nidumbur}
    \end{onecolentry}

    \vspace{0.10 cm}
    \begin{onecolentry}
        \begin{highlights}
            \item Machine learning-based prediction of cell-type-resolved brain eQTLs; co-author on the scEEMS preprint.
        \end{highlights}
    \end{onecolentry}

    \vspace{0.20 cm}

    \begin{onecolentry}
        \textbf{Clara Yu}, research intern, California Institute of Technology
    \end{onecolentry}

    \vspace{0.10 cm}
    \begin{onecolentry}
        \begin{highlights}
            \item Interpretable deep learning for functional fine-mapping of chromatin QTLs in Alzheimer's disease; first author of a peer-reviewed conference paper (IEEE ICBCB 2026).
        \end{highlights}
    \end{onecolentry}



\section{Fellowships, Grants, and Awards}

    \begin{twocolentry}{
        2026 – 2027
    }
        \textbf{Researcher}, NIH/NIA R01~AG086467, \textit{FunGen-xQTL: Unraveling the Genetic Basis of Molecular Functions in Alzheimer's Disease} (PI: Gao Wang, Columbia University)
    \end{twocolentry}

    \vspace{0.20 cm}

    \begin{twocolentry}{
        2018 – 2020
    }
        \textbf{NIH T32 Training Fellowship}, Harvard Medical School
    \end{twocolentry}

    \vspace{0.20 cm}

    \begin{twocolentry}{
        November 2018
    }
        \textbf{Stellar Abstract Award}, Harvard Program in Quantitative Genomics Conference
    \end{twocolentry}

    \vspace{0.20 cm}

    \begin{twocolentry}{
        May 2017
    }
        \textbf{Microsoft Azure Research Award} (\$100,000 in computing credits)
    \end{twocolentry}

    \vspace{0.20 cm}

    \begin{twocolentry}{
        June 2011
    }
        \textbf{AMS Mathematical Research Community}, Computational and Applied Topology
    \end{twocolentry}

    \vspace{0.20 cm}

    \begin{twocolentry}{
        2008, 2011
    }
        \textbf{NSF Conference Travel Grants} (MEGA 2011, Stockholm; Aspects of Moduli 2008, Pisa)
    \end{twocolentry}



\section{Technical Skills and Software}

    \begin{onecolentry}
        \textbf{Programming Languages:} Python, R, C/C++, Java, Stan
    \end{onecolentry}

    \vspace{0.2 cm}

    \begin{onecolentry}
        \textbf{Machine Learning:} PyTorch, TensorFlow, Pyro, scikit-learn, Hugging Face; experiment management with Optuna and Weights \& Biases
    \end{onecolentry}

    \vspace{0.2 cm}

    \begin{onecolentry}
        \textbf{Genomics:} genomic deep learning models (Enformer, ChromBPNet, DeepSEA); fine-mapping and heritability analysis (SuSiE, LDSC); whole-genome sequencing analysis at ADSP scale; single-cell genomics
    \end{onecolentry}

    \vspace{0.2 cm}

    \begin{onecolentry}
        \textbf{Infrastructure:} distributed computing (Spark, HPC/Slurm); cloud platforms (AWS, Azure, GCP); SQL and large-scale data engineering
    \end{onecolentry}

    \vspace{0.2 cm}

    \begin{onecolentry}
        \textbf{Released Software:} scEEMS, gruyere, parmigiano, LDmat (PyPI)
    \end{onecolentry}




\end{document}
